% Page 017

\seite{17}

\jahresueberschrift{Das Jahr 1889}
\pstart
Das Jahr 1889 brachte zunächst eine Veränderung in\ob
der Pfarrverwaltung. Am 27.\ Januar 1889 verstarb\ob
nach langer Krankheit der hochwürdige Herr Joh.\ Jocham
Kelz, Pfarrer in Seffern.

Geboren zu Fließem am 31.\ Januar 1831 empfing er\ob
am 16.\ Februar 1856 die Priesterweihe, wirkte als\ob
Kaplan zu C.\ Engen, dann in Losheim. 1861\ob
wurde er zum Pfarrer in Bettingen und 1868 zum\ob
Pfarrer in Seffern ernannt. Er war seinen Submittenten\ob
ein lieber Freund und seinen Pfarrkindern ein\ob
aufopfernder Seelsorger. Die Verwaltung der ver\ohb
waisten Pfarrei hatte nun zunächst Herr Pfarrer\ob
Bruck aus Bickendorf, bis am 8.\ Mai 1889\ob
der neuernannte Pfarrer Herr Michael Dunckel,\ob
bisher Kaplan in Mayen, die Pfarrstelle übernahm.

Am 15.\ Oktober erfolgte die Ernennung des Herrn\ob
Pfarrers Dunckel zum Lokalschulinspector für die\ob
Schulen Seffern, Sefferweich, Schleid und Heilenbach.\ob
Vor dem Herr Pfarrer Dunckel bekleidete Herr Bürger\ohb
meister Spang dieses Amt. etc.

Am 1.\ November 1889 wurde an die hiesige Schule der\ob
Herr Joh.\ Scholtes an Stelle des in Pension gegebenen\ob
Lehrer Dillenburg berufen. Derselbe war geboren zu\ob
Brück, im Kreise Wittlich und vorgebildet im
\pend
